\section*{Question 1}

Let $P=(p_{xy})_{x,y\in\{A,C,G,T\}}$ be the given $4\times4$ transition matrix
(rows and columns in the order $A,C,G,T$).
The root sequence is fixed as
\[
    X_0 = \mathrm{AAA},
\]
so at each site $k=1,2,3$ the root state is
\[
    R_k = A.
\]
The observed sequences are
\[
    X_1 = \mathrm{AGT},\quad
    X_2 = \mathrm{AAC},\quad
    X_3 = \mathrm{AGG}.
\]
Because sites evolve i.i.d., the likelihood of a tree topology $T$ is
\[
    L(T) = \prod_{k=1}^{3} L_k(T),
    \qquad
    \ell(T) = \log L(T) = \sum_{k=1}^{3} \log L_k(T),
\]
where $L_k(T)$ is the likelihood for site $k$.

Let $H_k\in\{A,C,G,T\}$ be the (unobserved) state at the internal node $h$
at site $k$.
Using the tree structures, the per-site likelihoods are:

\medskip\noindent
\textbf{Tree 1:} $0 \to h,\;h\to 1,\;h\to 2,\;0\to 3$
\[
    L_k(T_1)
    = p_{A,X_{3k}}\sum_{h} p_{A,h}\,p_{h,X_{1k}}\,p_{h,X_{2k}}.
\]

\noindent
\textbf{Tree 2:} $0 \to h,\;h\to 3,\;h\to 2,\;0\to 1$
\[
    L_k(T_2)
    = p_{A,X_{1k}}\sum_{h} p_{A,h}\,p_{h,X_{3k}}\,p_{h,X_{2k}}.
\]

\noindent
\textbf{Tree 3:} $0 \to h,\;h\to 1,\;h\to 3,\;0\to 2$
\[
    L_k(T_3)
    = p_{A,X_{2k}}\sum_{h} p_{A,h}\,p_{h,X_{1k}}\,p_{h,X_{3k}}.
\]

\medskip
Here $X_{ik}$ is the nucleotide at site $k$ in sequence $X_i$,
and $p_{xy}$ is the entry of $P$ for parent base $x$ and child base $y$.
The three site patterns (including the root) are
\[
    \begin{array}{c|cccc}
        k & R_k & X_{1k} & X_{2k} & X_{3k} \\ \hline
        1 & A   & A      & A      & A      \\
        2 & A   & G      & A      & G      \\
        3 & A   & T      & C      & G
    \end{array}
\]

Substituting these bases and the entries of $P$ into the formulas above
gives the following per-site likelihoods (numerical values):

\[
    \begin{array}{c|ccc}
            & L_1          & L_2                    & L_3                    \\ \hline
        T_1 & 0.9074011636 & 4.6821327\times10^{-5} & 1.7933930\times10^{-6} \\
        T_2 & 0.9074011636 & 4.6821327\times10^{-5} & 1.3051850\times10^{-6} \\
        T_3 & 0.9074011636 & 6.5951599\times10^{-3} & 1.3056000\times10^{-6}
    \end{array}
\]

Therefore,
\[
    \begin{aligned}
        L(T_1) & = L_1 L_2 L_3 \approx 7.62\times10^{-11},
               & \quad \ell(T_1)                           & \approx -23.30, \\
        L(T_2) & \approx 5.55\times10^{-11},
               & \quad \ell(T_2)                           & \approx -23.62, \\
        L(T_3) & \approx 7.81\times10^{-9},
               & \quad \ell(T_3)                           & \approx -18.67.
    \end{aligned}
\]

Since $T_3$ has the largest (least negative) log-likelihood,
the maximum-likelihood phylogenetic tree is

\[
    \boxed{\text{Tree 3: } 0 \to h,\; h \to 1,\; h \to 3,\; 0 \to 2.}
\]
